%% LyX 2.4.4 created this file.  For more info, see https://www.lyx.org/.
%% Do not edit unless you really know what you are doing.
\documentclass[english,t]{beamer}
\usepackage{mathptmx}
\usepackage[T1]{fontenc}
\usepackage[latin9]{inputenc}
\usepackage{amssymb}
\usepackage{cancel}
\PassOptionsToPackage{normalem}{ulem}
\usepackage{ulem}

\makeatletter
%%%%%%%%%%%%%%%%%%%%%%%%%%%%%% Textclass specific LaTeX commands.
% this default might be overridden by plain title style
\newcommand\makebeamertitle{\frame{\maketitle}}%
% (ERT) argument for the TOC
\AtBeginDocument{%
  \let\origtableofcontents=\tableofcontents
  \def\tableofcontents{\@ifnextchar[{\origtableofcontents}{\gobbletableofcontents}}
  \def\gobbletableofcontents#1{\origtableofcontents}
}

%%%%%%%%%%%%%%%%%%%%%%%%%%%%%% User specified LaTeX commands.
\usetheme{Warsaw}
% or ...

\setbeamercovered{transparent}
% or whatever (possibly just delete it)
\setbeamercovered{invisible} %makes overlays invisible

%gets rid of navigation symbols
\setbeamertemplate{navigation symbols}{}

%\usepackage{hyperref}
%\usepackage[usenames,dvipsnames,table]{xcolor} %adds additional color options
\definecolor{links}{HTML}{2A1B81}
%\hypersetup{colorlinks,linkcolor=blue,urlcolor=links}

\usepackage{multicol}

\usepackage{tikz}
\usetikzlibrary{calc}
\usetikzlibrary{plotmarks}
\usetikzlibrary{shapes}
\usepackage{pgfplots}
\pgfplotsset{compat=newest}

\usepackage{forest} %%for tree diagram

\usepackage{calc}
\usepackage[nomessages]{fp}

%\usepackage{advdate}    % Advancing/saving dates
\usepackage{datetime}   % Dates formatting
%\usepackage{datenumber} % Counters for dates

\usepackage[super]{nth}

\usepackage{etex}

\usepackage[outline]{contour}  %allows text halos
\contourlength{1pt} %sets halo width
%%example: \node at (0.75,0.5) {\contour{white}{\Large Text!}};
%%example:  \node [fill=white,inner sep=1pt] at (2.25,0.5) {\Large 			Text!};
%%example:  \node [fill=white,rounded corners=2pt,inner sep=1pt] at 			(3.75,0.5) {\Large Text!};

%%%%%%%%%%%%%%%%%%%%%%%%%%%
\def\formatdateny{\csname noyear\languagename\expandafter\endcsname\formatdate}
\def\noyearenglish#1, \the\@year{#1}
\let\noyearamerican\noyearenglish
\let\noyearbritish\noyearenglish

%%allows uncover with forest diagrams
\tikzset{
    invisible/.style={opacity=0,text opacity=0},
    visible on/.style={alt=#1{}{invisible}},
    alt/.code args={<#1>#2#3}{%
      \alt<#1>{\pgfkeysalso{#2}}{\pgfkeysalso{#3}} % \pgfkeysalso doesn't change the path
    },
}
\forestset{
  visible on/.style={
    for tree={
      /tikz/visible on={#1},
      edge={/tikz/visible on={#1}}}}}

%%%redefines column alignment to match vertically
%\define@key{beamerframe}{t}[true]{% top
  %\beamer@frametopskip=.2cm plus .5\paperheight\relax%
  %\beamer@framebottomskip=0pt plus 1fill\relax%
 % \beamer@frametopskipautobreak=\beamer@frametopskip\relax%
  %\beamer@framebottomskipautobreak=\beamer@framebottomskip\relax%
  %\def\beamer@initfirstlineunskip{}%
%}

%%provides pdf with 4 slides per page; don't forget to add 'handout' to remove overlays
%\usepackage{pgfpages}
%\pgfpagesuselayout{4 on 1}[a4paper, landscape, border shrink=7mm]
%\pgfpageslogicalpageoptions{1}{border code=\pgfusepath{stroke}}
%\pgfpageslogicalpageoptions{2}{border code=\pgfusepath{stroke}}
%\pgfpageslogicalpageoptions{3}{border code=\pgfusepath{stroke}}
%\pgfpageslogicalpageoptions{4}{border code=\pgfusepath{stroke}}
%%%Don't forget to add 'handout'

\makeatother

\usepackage{babel}
\begin{document}
\newcommand{\xmax}{2000}
\newcommand{\xmin}{0}
\newcommand{\xstep}{0} %must be xmin+n
\newcommand{\ymax}{500}
\newcommand{\ymin}{0}
\newcommand{\ystep}{0} %must be ymin+n

\newcounter{countEx} \setcounter{countEx}{0}
\newcounter{countProb} \setcounter{countProb}{0}
\resetcounteronoverlays{countEx} \resetcounteronoverlays{countProb}

\newcommand{\ex}[3]{
	\addtocounter{countEx}{1}
	\only<#1-#2>{\begin{exampleblock} {Example \chNum.\secNum.\arabic{countEx}.} #3	\end{exampleblock}}
}

\newcommand{\prob}[4]{
	\addtocounter{countProb}{1}
	\only<#1-#2>{\begin{exampleblock} {Problem  \chNum.\secNum.\arabic{countProb}.\,#3} #4	\end{exampleblock}}
}

\newcommand{\yline}[4]{
	(#2 - #4)/(#1 - #3)*(x - #1) + #2
}

\beamerdefaultoverlayspecification{<+->}

%%%%Chapter and Section numbers%%%%%
\newcommand{\chNum}{2}
\newcommand{\secNum}{4}
\newcommand{\secTitle}{Permutations and Combinations}

\title[Section \chNum.\secNum: \secTitle]{Section \chNum.\secNum: \secTitle}

\subtitle{Finite Mathematics~~MTH 107~~Fall 2014}

\author{Dr.\,James Ashe}

\titlegraphic{\includegraphics[height=2.75cm]{/home/jim/JamesAshe/MHU/images/MarsHilUniversityLogo-Virtical.png}}

\date{\formatdate{1}{9}{2014} and \formatdate{3}{9}{2014}}

\institute{Mars Hill University}

\makebeamertitle 
\begin{frame}{Permutations and Combinations}

\begin{block}{}
In section $2.3$, we selected \emph{one item each }from each category.
\end{block}
\begin{itemize}
\item [ex.]<2->Choose $1$ out of $6$ vegetables toppings and $1$ out
of $4$ meat toppings. How many different pizzas are possible?
\end{itemize}
\begin{block}<3->{}
In this section, we select 

\begin{itemize}
\item []<4->\emph{more than one item }from a category without replacement.
\item []<5->where \emph{order }does or does not matter,
\end{itemize}
\end{block}
\begin{itemize}
\item [ex.]<5->Choose $2$ out of $10$ people for president and vice president?
\item [ex.]<7->Choose $3$ out of $6$ vegetable toppings. How many different
pizzas are possible?
\end{itemize}
\end{frame}

\begin{frame}{Permutations}

\begin{itemize}
\item \ex{1}{}{$5$ students are giving presentations. How many different
lineups are possible?}

\begin{itemize}
\item []<2->$5\times4\times3\times2\times1=$ \only<3->{$=5!$}
\end{itemize}
\end{itemize}
\begin{block}<4->{When to use permutations?}

\begin{itemize}
\item <5->Items are selected from a single category.
\item <6->The \emph{order }is important
\item <7->Count using factorials, $n!$
\end{itemize}
\end{block}
\begin{itemize}
\item []<8->What if there is only time for $3$ out of the $5$ students
to give a presentation?
\end{itemize}
\end{frame}

\begin{frame}{Permutations}

\ex{1}{}{$3$ out of $5$ students are giving presentations. How
many different lineups are possible?}
\begin{itemize}
\item <2->There are no repeats.
\item []<3->$5$ 
\only<4->{$\times 4$} 
\only<5->{$\times 3$} 
\only<6->{$=60$}
\item []<7->Using factorials,
\item []<7->\only<8->{$\frac{5!}{2!}$} 
\only<9->{$=\frac{5\times4\times3\times2\times1}{2\times1}$} 
\only<10->{$=5\times4\times3=60$}
\end{itemize}
\begin{block}{Permutation Formula}
The number or \emph{permutations}, or arrangements, of $r$ items
selected \uline{without replacement} from a pool of $n$ items ($r\leq n$)
is,

\begin{itemize}
\item $_{n}P_{r}=\frac{n!}{(n-r)!}$
\end{itemize}
\end{block}
\end{frame}

\begin{frame}{Permutations}

\ex{1}{2}{A baseball team has $11$ players. How many different $9$ player lineups are possible?}
\begin{itemize}
\item []
\end{itemize}
\ex{2}{2}{At the end of the season, the team will award an "MVP", "most improved", and "LVP" awards. How many ways can the awards be given?}
\begin{itemize}
\item []
\end{itemize}
\ex{3}{}{The baseball league has $6$ teams. In how ways can they be ranked at the end of the season?}
\begin{itemize}
\item <4->[]Recall: $0!=1$
\item <5->No replacement, \uline{order matters}, and selected from $1$
category. 
\end{itemize}
\end{frame}

\begin{frame}{Combinations}

\begin{itemize}
\item <1->No replacement.
\item <2->Order matters \only<3->{$\Rightarrow$\emph{Permutations}, $_{n}P_{r}=\frac{n!}{(n-r)!}$}
\item []<4->What if order does not matter?

\begin{itemize}
\item [ex.]<5->How many $2$ element subsets can be selected from $A=\{a,b,c\}$?
\item []<6->$\{a,b\},\{a,c\},\{b,a\},\{b,c\},\{c,a\},\{c,b\}$

\begin{itemize}
\item []<7->$_{3}P_{2}=\frac{3!}{1!}=6$
\end{itemize}
\item []<8->Removing the repeats, $\frac{6}{2}=3$.
\item []<9->$\left(\frac{3!}{(3-2)!2!}\right)$
\end{itemize}
\end{itemize}
\end{frame}

\begin{frame}{Combinations}

\begin{block}{Combination Formula}
The number of distinct \emph{combinations} of $r$ items selected
without replacement from a pool of $n$ items $(r\leq n)$, 
\[
_{n}C_{r}=\frac{n!}{(n-r)!r!}
\]
Combinations are used whenever one or more items are selected (without
replacement) from a category and the order of selection is not important. 
\end{block}
\end{frame}

\begin{frame}{Combinations}

\ex{1}{}{A pizza place has $8$ toppings. If you choose $3$ different toppings, how many different pizzas are possible?}
\begin{itemize}
\item <2->[]Order does not matter.
\item <3->[]If order did matter, there would be 

\begin{itemize}
\item []<4->$_{8}P_{3}=\frac{8!}{(8-3)!}=\frac{8!}{5!}=$
\end{itemize}
\item <5->[]For each pizza with $3$ of the same toppings, there are $3!$
ways to place the order.

\begin{itemize}
\item [ex.]<6->$\{Pep_{1},Ham,Pep_{2}\}=\{Ham,Pep_{1},Pep_{2}\}=\{Ham,Pep_{2},Pep_{1}\},\dots$
\end{itemize}
\item <7->[]So, there are \only<8->{$\frac{8!}{(8-3)!}\frac{1}{3!}=56$
different pizzas.}
\end{itemize}
\end{frame}

\begin{frame}{Permutation or Combination?}

\prob{1}{1}{Find the indicated value.}{$11$ people try out for a baseball team, but there are only $9$ spots. How many different teams are possible?}

\prob{2}{2}{Find the indicated value.}{A committee of $3$ students must be selected for student council from a pool of $6$. How many different student councils are possible?}

\prob{3}{4}{Find the indicated value.}{A committee of $3$ students must be selected for student council from a pool of $6$. One will be president, one vice pressident, and one will be janitor. How many different student councils are possible?}

\prob{4}{4}{Find the indicated value.}{A committee of $3$ students must be selected for student council from a pool of $6$. One will be president and two will be vice pressidents. How many different student councils are possible?}

\prob{5}{5}{Find the indicated value}{A baseball league has $6$ teams. If every team must play every other team once in the first round of league play, how many games must be played?}
\end{frame}

\begin{frame}{When there is more than one category}

\ex{1}{}{A pool of $6$ seniors, $5$ juniors, $4$ sophomores, and
$3$ freshmen must select a committee of $4$. How many committees
are possible if the committee must contain
\begin{itemize}
\item [a.]<2->one person from each class
\item [b.]<3->any mixture of the classes
\item [c.]<4->$2$ seniors and $2$ juniors
\end{itemize}
}

\only<5-5>{We have $4$ categories\\
\hspace{2cm}$\underline{\,\,\,\,}\times\underline{\,\,\,\,}\times\underline{\,\,\,\,}\times\underline{\,\,\,\,}$}

\only<6-8>{No repeats and order does not matter}

\only<7-8>{ $\Rightarrow$ Combinations}

\only<8-8>{ \center $_{18}C_{4}=\frac{18!}{(18-4)!4!}$}

\only<9->{Now we have two categories}

\only<10-10>{\hspace{2cm} $\underline{\,\,\,\,\,\,\,}\times\underline{\,\,\,\,\,\,\,}$}

\only<11->{\hspace{2cm} $\underline{_{6}C_{2}}\times\underline{_{5}C_{2}}$}
\end{frame}

\begin{frame}{Distinguishable Permutations}

\ex{1}{1}{Find the number of \emph{distinguishable} permutations
of the letters of the work ``TRAIN''}

\ex{2}{}{Find the number of \emph{distinguishable} permutations of
the letters of the work ``TENNESSEE''}
\begin{itemize}
\item []<2->There are $9!$ ways to arrange $9$ items, but ``TENNESS\textcolor{blue}{E}\textcolor{red}{E}''
is not distinguishable from ``TENNESS\textcolor{red}{E}\textcolor{blue}{E}''.
\item <3->There are $4!$ ways to arrange $4$ ``E'''s,
\item <4->$2!$ ways to arrange $2$ ``N'''s,
\item <5->and $2!$ ways to arrange $2$ ``S'''s.
\item []<6->So there are $\frac{9!}{4!2!2!}=\frac{9\cdot8\cdot7\cdot6\cdot5}{4}=3,780$
\emph{distinguishable }arrangements.
\end{itemize}
\end{frame}

\begin{frame}{Problems}

\prob{1}{}{Refer to a standard 52 card deck (no jokers).}{How many five-card hands consisting of all hearts are possible?}
\begin{itemize}
\item []<2->One category, ``ways to get one suit'' $\Rightarrow\,\,\underline{\,\,\,\,\;}$.
\item []<3->No repeats; order does not matter $\Rightarrow$ combination
\item []<4->$_{13}C_{4}$\prob{5}{}{Refer to a standard 52 card deck (no jokers).}{How many five-card hands consisting of all the same suit are possible?}
\item []<6->No repeats; order does not matter $\Rightarrow$ combination
\item []<7->Two categories, ``ways to get a suit'' and ``ways to get
one suit'', $\underline{\,4\,}\times\underline{\,_{13}C_{4}\,}$.
\end{itemize}
\end{frame}

\begin{frame}{}

\prob{1}{}{Refer to a standard 52 card deck (no jokers).}{How many five-card hands consisting of 4 aces are possible?}
\begin{itemize}
\item []<2->Two categories: ``ways to get $4$ aces'' and ``ways to
get a non-ace'', $\underline{\;_{4}C_{4}\,}\times\underline{\;_{48}C_{1}\,}$.
\item []<3->No repeats; order does not matter $\Rightarrow$ combination
\end{itemize}
\prob{4}{}{Refer to a standard 52 card deck (no jokers).}{How many five-card hands consisting of 3 kings and 2 queens are possible?}
\begin{itemize}
\item []<5->Two categories: ``ways to get $3$ kings'' and ``ways to
get $2$ queens, $\underline{\;_{4}C_{3}\,}\times\underline{\;_{4}C_{2}\,}$.
\item []<6->No repeats; order does not matter $\Rightarrow$ combination
\end{itemize}
\end{frame}

\end{document}
